% Compile with: latexmk -pdf paper.tex (from docs/paper, where refs.bib and figs/ live)
\documentclass[10pt]{article}
\usepackage[margin=1in]{geometry}
\usepackage{times}
\usepackage{amsmath,amssymb,amsthm}
\usepackage{graphicx}
\usepackage{booktabs}
\usepackage{microtype}
\usepackage{needspace}
\usepackage[authoryear,round]{natbib}
\usepackage{hyperref}

\newtheorem{theorem}{Theorem}
\newtheorem{proposition}[theorem]{Proposition}
\newtheorem{corollary}[theorem]{Corollary}
\newtheorem{definition}{Definition}
\newtheorem{lemma}{Lemma}[section]
\newcommand{\ess}{\mathrm{ESS}}
\newcommand{\E}{\mathbb{E}}
\newcommand{\R}{\mathbb{R}}

\title{How Many Demonstrations Is a Task Description Worth?}
\author{Bruce Quan Nguyen}
\date{\today}

\begin{document}

\maketitle

\begin{abstract}
Prompts for in-context learning usually combine a task description with demonstrations, yet theories of in-context learning mostly model only the demonstrations. We ask how many demonstrations a description is worth. We define the \emph{effective sample size} (ESS) of a description as the number of demonstrations that lower a Bayes-optimal predictor's regret by the same amount, and study it in in-context linear regression. The ESS is set by the description's \emph{specificity}, how much uncertainty about the task it removes, and its \emph{reliability}, the probability that it is valid. For a reliable description, we prove that the ESS grows linearly in the task dimension when the description is loose and in inverse proportion to the variance it leaves when it is specific. For a description that may be invalid, the ESS without demonstrations stays bounded however specific the description is, and a predictor that trusts it fully can do worse than one that ignores it. Meta-trained Transformers come close to the Bayes-optimal ESS for reliable descriptions and inherit their trust in a description from their training distribution. Pretrained language models that read the description in words respond only weakly to its stated precision: across six OLMo~2 models and three wordings, the weight they give a description depends on the model and the wording far more than on the precision it states.
\end{abstract}

\section{Introduction}
\begin{figure}[t]
\centering
\includegraphics[width=\linewidth]{figs/overview.pdf}
\caption{(a) A prompt read as Bayesian inference, and the ESS of its description. (b) The prior on the task given three descriptions (sketched, not to scale), and the ESS of each with no demonstrations in hand, at $d = 16$ and $a_0 = 10$. Specificity raises the ESS (\S\ref{sec:specificity}); a one-in-ten chance of being invalid removes most of that gain (\S\ref{sec:reliability}).}
\label{fig:overview}
\end{figure}
A prompt for in-context learning (ICL) usually carries two kinds of information about the task: a description in natural language and a few demonstrations \citep{brown2020}. In practice the description is an instruction; we read it as a statement about the task, since its value lies in what it says the task is, and set aside how a model is made to follow it. The evidence on what a description contributes is mixed. In prompt-based fine-tuning a well-crafted prompt can be worth hundreds of data points \citep{lescao2021}, and a written summary can stand in for many demonstrations \citep{honda2025}, yet models also perform adequately with misleading descriptions \citep{webson2022} and can prioritize demonstrations over explicit instructions \citep{gupta2025}. Theories of ICL mostly model the prompt as a sequence of demonstrations \citep{garg2022, akyurek2023, xie2022}. Those that add a description show that a learner can use it to predict better \citep{huangge2025}, but not how much it is worth in the unit a prompt writer trades it against: demonstrations. We ask the question directly: how many demonstrations is a task description worth, and what determines the number?

We take the Bayesian view of ICL \citep{xie2022}. Pretraining gives the model a prior over a latent task, and the prompt is evidence about that task. A demonstration is one noisy observation of the task's input-output behaviour. A description is a statement about the task itself, and conditioning on it gives a sharper prior. Bayesian statistics measures the information in a prior by its effective sample size (ESS), the number of observations that would carry the same information \citep{morita2008}. We define the ESS of a description in the same way and compute it in noisy in-context linear regression with a Gaussian task prior \citep{garg2022, akyurek2023, raventos2023}, where Bayes-optimal prediction has a closed form.

Two properties of a description determine its ESS. The first is its \emph{specificity}: a description may state exactly how to do the task or only loosely indicate it, and the more specific it is, the less uncertainty about the task it leaves. The second is its \emph{reliability}: the probability that it is valid, since a description can conflict with the task the demonstrations come from \citep{evans2006}. Our contributions are:
\begin{itemize}
    \item \textbf{A common unit.} We define the ESS of a task description as the number of demonstrations that lower a Bayes-optimal predictor's regret by the same amount (Definition~\ref{def:ess}). It puts descriptions and demonstrations on one scale.
    \item \textbf{What specificity is worth.} For a reliable description, we prove that the ESS grows linearly in the task dimension $d$ when the description is loose and linearly in $1/r$ when it is specific, where $r$ is the fraction of the prior variance the description leaves (Theorem~\ref{prop:reliable}). Roughly, it is $(d-1)(1-r) + 1/(r a_0)$.
    \item \textbf{What unreliability costs.} For a description that is valid with probability $p$, the regret is that of a predictor told whether the description is valid, plus the information the next label carries about validity, which totals at most the entropy $H(p)$ over the whole prompt (Proposition~\ref{prop:unreliable}). Consequently, the ESS without demonstrations is bounded however specific the description is, and with many demonstrations it is at most $p$ times that of a reliable description (Corollary~\ref{cor:ess}). A predictor that instead trusts the description fully pays a penalty that grows with specificity, and a specific description that is sometimes invalid can then be worse than none.
\end{itemize}
We then measure the ESS in two kinds of learner. Meta-trained Transformers come close to the Bayes-optimal ESS for reliable descriptions, and their trust in a description follows the reliability they were trained on (\S\ref{sec:networks}). Pretrained language models, given a description in words, give it a weight that depends on the model and the wording but only weakly on the precision it states (\S\ref{sec:llm}).

\section{Related Work}
\paragraph{ICL as Bayesian inference.} Our approach builds on the understanding of ICL as implicit Bayesian inference \citep{xie2022}, in which meta-trained Transformers have been shown to implement optimal estimators such as Bayes-optimal ridge regression \citep{akyurek2023, raventos2023, panwar2024}.

\paragraph{Task descriptions in theories of ICL.} Some of this literature includes task descriptions, modeling them as tokens that indicate hypothesis classes or carry input means \citep{huangge2025, lin2025, tong2026, lin2024}, but these works generally do not quantify the description's value in units of examples, nor do they formalize the probability that a description is invalid. Closest to our setting, \citet{zhu2026} show that a Transformer given a prefix of related datasets performs amortised hierarchical Bayesian prediction; we express the value of such conditioning information in demonstrations and allow it to be invalid.

\paragraph{Effective sample size of a prior.} Our ESS follows the prior sample size methods of Bayesian statistics \citep{morita2008, reimherr2021, neuenschwander2020}, particularly those that use predictive regret \citep{reznik2026}. To model a description that may be invalid we use a robust mixture prior, a technique from clinical trials for historical data that may conflict with the trial \citep{schmidli2014}.

\paragraph{Descriptions and priors in pretrained language models.} Empirical work has priced prompts in training examples: \citet{lescao2021} in prompt-based fine-tuning, and \citet{cook2026}, who find that one instruction in fine-tuning data can stand in for up to a hundred execution examples. \citet{wenliang2025} treat a prompt as conditioning for a Bayesian sequence predictor and find that an optimized prompt can be shorter than demonstrations of equal effect. Tests of pretrained models against a known optimum on synthetic tasks find that in-context evidence outweighs a stated prior \citep{gupta2025}, that learning curves approach a Bayesian reference \citep{akata2026}, and that models estimate a density from numbers in context like a kernel estimator of adaptive width \citep{liu2025}. We add to such a test a description of stated precision that may be invalid, and price it in demonstrations against the Bayes-optimal value.

\section{Problem Setup}\label{sec:setting}

\paragraph{Tasks and demonstrations.} A task is a weight vector $w \in \R^d$. A demonstration is an input $x \sim \mathcal{N}(0, I_d)$ with a label $y = w^\top x + \epsilon$, where $\epsilon \sim \mathcal{N}(0, \sigma_y^2)$. The base prior over tasks is $w \sim \mathcal{N}(0, s_0^2 I_d)$. The prior signal-to-noise ratio is $a_0 = s_0^2/\sigma_y^2$; we set $\sigma_y = 1$ throughout and $a_0 = 10$ in the numerical results.

\paragraph{Descriptions.} A description is a vector $m \in \R^d$. A latent indicator $Z \sim \mathrm{Bernoulli}(p)$ says whether the description is \emph{valid} ($Z = 1$) or \emph{invalid} ($Z = 0$):
\begin{align}
    m &\sim \mathcal{N}(0, (1-r)s_0^2 I_d) \\
    w \mid m, Z=1 &\sim \mathcal{N}(m, r s_0^2 I_d) \\
    w \mid m, Z=0 &\sim \mathcal{N}(0, s_0^2 I_d)
\end{align}
A valid description locates the task up to a variance $r s_0^2$ in every direction; an invalid one is independent of the task. Marginalizing over $Z$, the prior for $w$ given $m$ is the robust mixture $p\,\mathcal{N}(m, rs_0^2 I_d) + (1-p)\,\mathcal{N}(0, s_0^2 I_d)$ \citep{schmidli2014}. The parameter $r \in (0, 1)$ is the fraction of the prior variance that a valid description leaves, and it sets the description's \textbf{specificity}: the smaller $r$, the more specific the description, from one that only loosely indicates the task ($r$ near 1) to one that all but states it ($r$ near 0). The parameter $p$ is the description's \textbf{reliability}. We call a description with $p = 1$ \emph{reliable}.

\paragraph{Regret and ESS.} A predictor sees a prompt $D$ (the description, $n$ demonstrations, and a query $x$) and gives a predictive density $f(y \mid D)$ for the query's label. We measure it by its expected one-step log-loss regret against an oracle that knows the task $w$, the per-step regret of \citet{genewein2025}:
\begin{equation}\label{eq:regret}
    R = \E \left[ -\log f(y \mid D) + \log \phi(y; w^\top x, \sigma_y^2) \right],
\end{equation}
where $\phi$ is the Gaussian density. Let $R^{\mathrm{desc}}(n)$ be the Bayes-optimal regret with the description and $n$ demonstrations, $R^{\mathrm{ex}}(n)$ the Bayes-optimal regret with $n$ demonstrations alone, and $R^{\mathrm{desc}}_{p=1}(n)$ the former for a reliable description.
\begin{definition}\label{def:ess}
The effective sample size (ESS) of a description, given $n$ demonstrations, is $\ess_n = n^* - n$, where $n^*$ is the root of $R^{\mathrm{ex}}(n^*) = R^{\mathrm{desc}}(n)$, with $R^{\mathrm{ex}}$ linearly interpolated between integers. The standalone ESS is $\ess_0$, which we also write $\ess$.
\end{definition}
The ESS is the number of further demonstrations a predictor without the description would need to match the predictor with it. It is the data-dependent prior sample size of \citet{reimherr2021} with predictive regret as the measure of uncertainty. For the Bayes-optimal predictor it is never negative. The same definition applies to any other learner through its own two regret curves; there it is negative when the description hurts, and when the regret with the description is at least the regret with nothing there is no root and we write $\ess_n \le -n$.

\section{Specificity: Reliable Descriptions}\label{sec:specificity}
With a reliable description ($p = 1$) the posterior over $w$ is Gaussian with some covariance $\Sigma$, and the regret is half the expected log of the predictive variance relative to the noise, $R = \tfrac12 \E \log(1 + x^\top \Sigma x)$.

\begin{theorem}[ESS of a reliable description]\label{prop:reliable}
Let $p = 1$ and let $c = 1/(r a_0) = \sigma_y^2/(r s_0^2)$ be the conjugate-prior ESS.
\begin{enumerate}
    \item[(i)] (Loose descriptions) If $c \le 1$, then as $rd \to \infty$, $\ess = (d-1)(1-r) + (1-r)c + O(1/(rd))$.
    \item[(ii)] (Specific descriptions) If $c \ge d$, then as $c/d \to \infty$, $\ess = c + d + 1 - 1/a_0 + O(d/c)$.
    \item[(iii)] For every $d$, $a_0$, and $r$, $\ess \le c + d + 2$.
\end{enumerate}
\end{theorem}
Appendix~\ref{app:specificity} states the theorem in a fuller form and proves it. In both regimes the ESS is the sum of two terms. The first counts dimensions: a description that removes a fraction $1-r$ of the prior variance replaces about $(1-r)(d-1)$ demonstrations, because at high signal-to-noise ratio each of the first $d$ demonstrations removes the uncertainty in one direction. The second is the conjugate-prior ESS $c$, the number of observations with noise $\sigma_y^2$ that carry as much information as a prior of variance $r s_0^2$ \citep{neuenschwander2020}. A single formula summarizes both regimes: $\ess \approx (d-1)(1-r) + 1/(ra_0)$. It differs from the expansion in (i) by $1/a_0$ and from the one in (ii) by less than two demonstrations. It is not a theorem for $1 < c < d$, where only the bound (iii) applies, but at every $r$ in Figure~\ref{fig:ess}a it is within two demonstrations of the exact ESS. At $d = 16$ and $a_0 = 10$, a description that removes nine tenths of the prior variance ($r = 0.1$) is worth $14.9$ demonstrations, and one that leaves a thousandth ($r = 0.001$) is worth $117$ (Figure~\ref{fig:ess}a). With many demonstrations in hand the dimension term fades, since the demonstrations have already removed that uncertainty (Corollary~\ref{cor:ess}): at $r = 0.1$ the ESS falls from $14.9$ alone to $1.1$ with $100$ demonstrations in hand.

\begin{figure}[t]
\centering
\includegraphics[width=\linewidth]{figs/ess.pdf}
\caption{ESS of a description against $r$, the fraction of the prior variance it leaves, with $n$ demonstrations in hand ($\ess_n$ of Definition~\ref{def:ess}); one panel per reliability $p$. The ESS of a reliable description grows without bound as it becomes more specific (a). When the description may be invalid, the ESS with no demonstrations is capped at the dotted limit, and many demonstrations in hand lift it (b, c). $d = 16$, $a_0 = 10$, eight batches of 8{,}000 prompts, with standard errors below $0.1$ demonstrations ($0.2$ for the dotted line at $p = 0.99$). Dotted: the limit of the standalone ESS as $r \to 0$.}
\label{fig:ess}
\end{figure}

\section{Reliability: Descriptions That May Be Invalid}\label{sec:reliability}
When $p < 1$ the Bayes-optimal predictive is a mixture of two predictives, one that takes the description as valid and one that ignores it, weighted by the posterior probability that the description is valid.

\begin{proposition}[Regret of a description that may be invalid]\label{prop:unreliable}
For every $n \ge 0$, with $D_n$ the prompt with $n$ demonstrations and $y$ the label of its query,
\begin{equation}\label{eq:decomp}
    R^{\mathrm{desc}}(n) \;=\; p\, R^{\mathrm{desc}}_{p=1}(n) + (1-p)\, R^{\mathrm{ex}}(n) + I_n, \qquad I_n = I(Z;\, y \mid D_n),
\end{equation}
where $I_n$ is the mutual information between the validity $Z$ and the label given the prompt. Moreover $I_n = H(Z \mid D_n) - H(Z \mid D_{n+1})$, so $0 \le I_n \le H(Z \mid D_n) \le H(p)$ and $\sum_{n \ge 0} I_n \le H(p)$, with $H(p)$ the binary entropy.
\end{proposition}
Appendix~\ref{app:reliability} gives the proof. The first two terms are the regret of a predictor that is told whether the description is valid: with probability $p$ it holds a reliable description, and with probability $1-p$ it holds only the demonstrations. So a description that may be invalid is worth what a reliable one is worth when it is valid and nothing when it is not.

The last term is the price of not being told, and under log loss it is exactly what the label reveals about validity. It depends on the setting, through $d$, $a_0$, $r$, and $n$: it vanishes when the two predictives agree, so that the label says nothing about validity, and it is largest when a specific description makes them differ sharply ($0.04$ nats at $r = 0.5$ and $0.26$ at $r = 0.001$, with no demonstrations at $d = 16$, $a_0 = 10$, $p = 0.9$). But there are only $H(p)$ nats to learn about $Z$, so the price summed over the whole prompt is at most $H(p)$, whatever $d$, $a_0$, and $r$: $0.33$ nats at $p = 0.9$. This is an average over valid and invalid descriptions. Case by case, the price summed over the whole prompt is at most $\log(1/p)$ when the description is valid and $\log(1/(1-p))$ when it is invalid, $0.11$ and $2.3$ nats at $p = 0.9$, because the mixture starts with weight $p$ or $1-p$ on the right predictive; the invalid case is costly but rare. The decomposition is the standard one for a mixture forecaster under log loss \citep{cesabianchi2006, merhav1998}; what is new is what it implies for the ESS.

\needspace{7\baselineskip}
\begin{corollary}[ESS of a description that may be invalid]\label{cor:ess}
Fix $d$, $a_0$, and $r$.
\begin{enumerate}
    \item[(i)] (Without demonstrations) If $p < 1$, then $\ess_0 \le n^*_p$, where $n^*_p$ is the finite root of $R^{\mathrm{ex}}(n^*_p) = (1-p)\,R^{\mathrm{ex}}(0)$, which does not depend on $r$.
    \item[(ii)] (With many demonstrations) If $p = 1$, then $\ess_n \to c - 1/a_0$ as $n \to \infty$. If $p < 1$, then $\limsup_{n \to \infty} \ess_n \le p\,(c - 1/a_0)$.
\end{enumerate}
\end{corollary}

\paragraph{Without demonstrations the ESS is bounded.} The average regret cannot fall below the share contributed by invalid descriptions, $(1-p)R^{\mathrm{ex}}(0)$, however specific the description is, and (i) turns that floor into a bound on the ESS. At $d = 16$, $a_0 = 10$, and $p = 0.9$ the bound $n^*_p$ is $40$ demonstrations. It applies even to the predictor told whether the description is valid, so it comes from invalid descriptions being worth nothing, not from doubt; not being told lowers the limit of the exact ESS as $r \to 0$ to $25$ (dotted in Figure~\ref{fig:ess}c). Refining $r$ from $0.1$ to $0.001$ raises the ESS of a reliable description from $14.9$ to $117$, but that of a $p = 0.9$ description only from $13.1$ to $23.0$.

\paragraph{With many demonstrations the ESS is at most $p$ times the reliable one.} Demonstrations shrink the regret in the invalid case, $R^{\mathrm{ex}}(n)$, and so lift the bound. In the limit a reliable description is worth $c - 1/a_0$ demonstrations: the dimension term is gone and what remains is the precision the description adds to the prior. A description that may be invalid is then worth at most $p$ times as much. With $100$ demonstrations in hand, the $r = 0.001$ description is worth $109$ demonstrations if reliable and $87.4$ at $p = 0.9$ (Figure~\ref{fig:ess}), against limits of $99.9$ and at most $89.9$.

\paragraph{Full trust.} The Bayes-optimal predictor hedges, and Proposition~\ref{prop:unreliable} prices the hedge at no more than $H(p)$ on average. A predictor that instead treats the description as certainly valid predicts as if $p = 1$ whether or not the description is valid. With no demonstrations its regret is
\begin{equation}\label{eq:trust}
    R^{\mathrm{trust}}(0) \;=\; R^{\mathrm{desc}}_{p=1}(0) \;+\; (1-p)(1-r)\,\E\!\left[\frac{a_0\|x\|^2}{1 + r a_0 \|x\|^2}\right]
\end{equation}
(Appendix~\ref{app:trust}): the regret of a reliable description plus a penalty for the invalid case. The penalty grows as the description becomes more specific, up to $(1-p)\,a_0 d$ as $r \to 0$, and $H(p)$ does not limit it. When the penalty exceeds what the description gains, full trust is worse than ignoring the description. At $d = 16$, $a_0 = 10$, and $p = 0.9$, the $r = 0.001$ description costs the fully trusting predictor $13.7$ nats, against $2.51$ with no description; at $p = 0.99$ the same description costs it $1.43$ nats and still helps.

\section{Meta-Trained Transformers}\label{sec:networks}
The Bayes-optimal predictor knows how specific and how reliable descriptions are. A network can learn both from the distribution it is trained on. We therefore meta-train Transformers on the setting of \S\ref{sec:setting} and ask how close they come to the Bayes-optimal ESS, and where their trust in a description comes from.

\paragraph{Setup.} Prompts use the prefix embedding of \citet{huangge2025}: an optional description token (their descriptor token) carrying $m$, up to $20$ demonstration tokens each holding an input beside its label, and a query token, with two indicator coordinates marking the token type and no positional encoding. The architecture and optimiser follow \citet{garg2022}: $12$ layers, $8$ heads, width $256$, Adam at a learning rate of $10^{-4}$, fresh prompts at every step, $40{,}000$ steps at batch $1024$ without a curriculum, at $d = 5$ and $a_0 = 10$. The network outputs a mean and a log variance for the query's label and is trained with log loss, so its regret is on the scale of~\eqref{eq:regret}. The number of demonstrations is uniform on $0$ to $20$ and the description is omitted in half the prompts. We train one network for each combination of a specific ($r = 0.01$) or loose ($r = 0.1$) description with reliability $p \in \{1, 0.99, 0.9\}$, one training run per setting and a second seed at $p = 0.9$, $r = 0.01$. The description token supplies only $m$, so each network must learn $r$ and $p$ from its training distribution. We evaluate each network on $20{,}000$ fresh prompts for every $n$ from $0$ to $20$, with and without a description, on prompts of all three reliabilities. Its ESS follows Definition~\ref{def:ess} with its own regret curves. We compare it, on the same prompts, with the Bayes-optimal predictor that uses the reliability the network was trained on (\emph{trained}) and the one that uses the reliability of the test prompts (\emph{calibrated}) (Table~\ref{tab:networks}). A learner's \emph{implied trust} is the weight $q$ for which the predictor that uses the mixture prior of \S\ref{sec:setting} with $q$ in place of $p$ best matches the learner's mean predictions.

\begin{table}[t]
\centering\small
\caption{Networks against the reference predictors ($d=5$, $a_0=10$; 20{,}000 evaluation prompts per network, with the standard error of its ESS; the last three rows are the second seed). ``Trained'' and ``calibrated'' are the Bayes-optimal predictors with the training and the test reliability, computed once on 200{,}000 prompts (standard errors below $0.12$) and shown once where they coincide; ``single Gaussian'' is the best predictive that a mean and a variance can express, where the Bayes-optimal predictive is a mixture. ``Trust'' is the network's implied trust with the description alone, on a grid of step $0.05$ plus $0.99$ and $0.999$. The last two columns are the mean regret gap, network minus the trained predictor, over $n = 0, \dots, 20$, with and without a description.}
\label{tab:networks}
\vspace{3pt}
\input{../../results/rq3-meta-trained/networks_table.tex}
\end{table}

\paragraph{Networks approach the Bayes-optimal ESS within what their output can express.} Trained and tested on reliable descriptions, the networks come close to the Bayes-optimal ESS: $14.7 \pm 0.4$ against $15.6$ at $r = 0.01$ and $5.24 \pm 0.05$ against $5.25$ at $r = 0.1$. With descriptions that may be invalid and no demonstrations, they fall short of it: at $p = 0.9$, $4.30 \pm 0.07$ against $6.92$ and $3.57 \pm 0.06$ against $4.26$, and at $p = 0.99$, $8.9 \pm 0.2$ against $12.9$ and $5.01 \pm 0.05$ against $5.08$. The cause is the output. The Bayes-optimal predictive is then a mixture of two Gaussians, which a network that emits one mean and one variance cannot express. The networks match the best single Gaussian instead, whose ESS is $4.29$, $3.62$, $8.54$, and $4.93$ in the same four cases. As demonstrations settle whether the description is valid, the mixture collapses to one component and the networks rejoin the Bayes-optimal predictor: at $r = 0.01$, $p = 0.9$ the network's regret is within $0.05$ nats of it from $n = 6$ on, and its ESS with ten demonstrations in hand is $8.1$ against $9.0$. A second seed reproduces the first ($4.26 \pm 0.08$ alone, $8.3$ with ten demonstrations). The shortfall is not under-training: a network trained at $r = 0.05$, $p = 0.9$ and continued to $80{,}000$ steps kept it ($3.9$ against $5.0$).

\paragraph{Trust is inherited from training.} A network's implied trust follows the reliability it was trained on. Networks trained only on reliable descriptions ($p = 1$) have an implied trust near $1$ and behave as the fully trusting predictor of \S\ref{sec:reliability}. When one description in ten is invalid, the network trained on specific descriptions has a regret of $2.87$ nats with the description alone against $1.87$ without it ($3.32$ for the fully trusting predictor, which the network beats by hedging slightly, at an implied trust of $0.999$), so its ESS is at most zero where the calibrated predictor obtains $6.9$. Networks trained at $p = 0.99$ have an implied trust of $0.95$ and lose the specific description's entire value when one in ten is invalid ($0.1 \pm 0.4$ against $6.9$). Networks trained at $p = 0.9$ keep hedging, with an implied trust of $0.85$ to $0.9$ across the two seeds, even on prompts whose descriptions are all valid, where their ESS is $6.4$ and $6.6$ against the calibrated $15.6$ at $r = 0.01$.

\section{Pretrained Language Models}\label{sec:llm}
The networks of \S\ref{sec:networks} know how specific and how reliable descriptions are because they were trained on the task distribution. A pretrained language model has had no such training: it has to take a description as it is written. We measure what a description is worth to such a model, with the Bayes-optimal predictor as the yardstick. The test does not assume that language models are Bayesian. Its design follows earlier tests of pretrained language models against a known optimum on synthetic tasks \citep{gupta2025, akata2026}.

\paragraph{Setting.} We use the simplest case of \S\ref{sec:setting}: one weight ($d = 1$) and the input fixed at $x = 1$, so that a task is a hidden mean $w$ and a demonstration is a number $y = w + \epsilon$. With inputs, a line about the weights must also say how the inputs combine, which is information that the description of \S\ref{sec:setting} does not carry, and using it takes arithmetic. With a single hidden mean, the description is the only instruction in the prompt. In this case the Bayes-optimal regret with prior precision $\tau$ and $n$ numbers is $\tfrac12 \log(1 + 1/(\tau + n))$, which depends on $\tau + n$ alone, so a reliable description is worth $\ess_n = c - 1/a_0$ numbers at every $n$, up to the interpolation in Definition~\ref{def:ess}; this is the limit in Corollary~\ref{cor:ess}(ii). The case has no dimension term, so this test does not cover that part of Theorem~\ref{prop:reliable}, which the networks of \S\ref{sec:networks} do.

\Needspace{9\baselineskip}
\paragraph{Prompts.} A prompt is a list of numbers, one per line, preceded in the description conditions by one line:

\smallskip\noindent\begin{minipage}{\linewidth}
\begin{verbatim}
    Mean: 508 ± 10
    516
    497
    523
    ...
\end{verbatim}
\end{minipage}
\smallskip

\noindent There is no preamble, no instruction about the form of the answer, and no chat template; the model continues plain text, as in \citet{liu2025}, whose prompts are numbers alone. The hidden mean is drawn from $\mathcal{N}(500, 100^2)$ and the numbers have noise of standard deviation $30$ and are rounded to integers (a shift of the prior mean leaves the Bayes-optimal results unchanged). Every number is then a positive three-digit integer (a task with a value outside $100$ to $999$ is redrawn), because language models predict numbers of few digits best \citep{gruver2023}, and wider ranges and negative values degrade their predictions \citep{requeima2024}. The line states the description's mean $m$ and its tolerance $\sqrt{r}\,s_0$, which the Bayes-optimal predictor of \S\ref{sec:setting} knows and the networks of \S\ref{sec:networks} learned in training. The tolerance is $30$, $10$, or $3$, for which $c = 1$, $9$, and $100$ ($r \approx 0.1$, $0.01$, and $0.001$ at $a_0 \approx 11$, the range of Figure~\ref{fig:ess}). An \emph{exact} line states the mean with no tolerance, and an \emph{invalid} line states the value of an independently drawn task. To see how much the wording matters, we repeat the valid and invalid conditions with two other wordings: ``The mean is about 508, give or take 10.'' and ``Mean in [488, 528]'', an interval of two tolerances either side.

\paragraph{Models and measurement.} We use OLMo~2 \citep{walsh2025} at $7$, $13$, and $32$ billion parameters, each as the base model and as its instruction-tuned version. Its weights are openly licensed; base and instruction-tuned models differ in how much weight they give to stated information \citep{gupta2025}; and its tokenizer encodes every three-digit integer as one token. One pass over a prompt of $200$ numbers, the context length of \citet{liu2025}, therefore gives the model's next-token distribution before every number \citep{kossen2024}. Restricted to the $900$ number tokens and renormalized, it is the model's predictive distribution for the next number, given that a number comes next, which we read as \citet{liu2025} and \citet{requeima2024} do. The regret~\eqref{eq:regret} at each position is the expected log loss of that distribution under the true distribution of the number, minus the oracle's; taking the expectation exactly, instead of scoring the one observed number, removes the sampling noise of the label. The model's ESS follows Definition~\ref{def:ess} with its own two mean regret curves, with and without the line, the latter made monotone by isotonic regression. We average $\ess_n$ over $n = 5, \dots, 10$, a window fixed before the runs, because the models are not told the noise level and do not expect a number at all before the first few. Every condition uses the same $2{,}000$ tasks, so conditions differ only in the line; in a simulation of the ideal learner on this design, $2{,}000$ tasks recover its ESS as $0.94 \pm 0.11$, $8.9 \pm 0.5$, and $100 \pm 5$ at the three tolerances (mean and standard deviation over $100$ replications). Intervals are $95\%$ bootstrap intervals over tasks.

A line can change a model's predictions whatever it says, so we add two controls. The \emph{content ESS} of a valid line is Definition~\ref{def:ess} with the invalid line of the same tolerance as the reference in place of the prompt with no line: the number of extra numbers that a model holding the wrong line needs to match the model holding the right one. It holds the form of the prompt fixed, and it measures what the content of the line adds: the gain from the right value together with the harm from the wrong one, for a model that does not discard it. It is not a bound on the ESS against no line, which also includes the effect of the form. And we repeat every ESS with the full next-token loss, which also charges the model for the probability it puts on tokens other than numbers.

We also read the mean of the predictive distribution. For the ideal learner with a valid line it is $(c\,m + n\,\bar y)/(c + n)$, a weighted average of the stated value and the mean $\bar y$ of the numbers so far. At each $n$ we regress the model's predictive mean on $m$ and $\bar y$ across tasks, with a constant, and call $n\,a_n/b_n$, the ratio of the two coefficients scaled by $n$, the \emph{implied strength} of the line: the number of numbers the model weighs it as. It is the conjugate-prior ESS $c$ of \S\ref{sec:specificity} for the ideal learner with a valid line, and it does not depend on how well the model learns the spread of the numbers. We average it over the same window, and we also record the probability the model puts on exactly the stated number.

\begin{table}[t]
\centering\small
\setlength{\tabcolsep}{3.5pt}
\caption{A one-line description with a stated tolerance of $\pm 10$, for OLMo~2 models and three wordings ($\pm$: ``Mean: $m \pm 10$''; sent.: the sentence; int.: the interval); $2{,}000$ tasks, $n = 5, \dots, 10$ numbers in hand. ``Strength'' is the implied strength of the line, for a valid and an invalid line. ``Content ESS'' is the ESS of a valid line measured against the invalid line of the same wording, and ``ESS vs.\ no line'' is Definition~\ref{def:ess} as stated, both over the number tokens. The last two columns are the regret with no line. The first Bayes-optimal row is the predictor that takes the line as valid; the second is the predictor that hedges, with a prior probability of validity from $0.5$ to $0.9$. Bootstrap $95\%$ intervals are within $\pm 0.5$ of each ESS and within $\pm 1.0$ of each strength, except the 32B instruction-tuned model's sentence ($\pm 2.7$). Table~\ref{tab:llm-full} has every tolerance.}
\label{tab:llm}
\vspace{3pt}
\input{../../results/llm-test/llm_table.tex}
\end{table}

\begin{figure}[t]
\centering
\includegraphics[width=\linewidth]{../../results/llm-test/llm_test.pdf}
\caption{Implied strength of a valid line against its stated tolerance, one panel per wording, for every model (dashed: instruction-tuned). Dotted: the Bayes-optimal strength $c$, which rises from $1$ to $100$.}
\label{fig:llm}
\end{figure}

\paragraph{The weight of a description depends on the model and the wording, and only weakly on the precision it states.} At $\pm 10$, where the Bayes-optimal predictor that takes the line as valid weighs the stated value as $9$ numbers and one that hedges weighs it as $4.4$ to $7.5$, the models weigh it as $0.6$ to $18$ (Table~\ref{tab:llm}). Instruction-tuned models give it more weight than base models of the same size, the 32B model more than the smaller ones of its kind, and, for five of the six models, the sentence more than the other wordings: the 32B instruction-tuned model weighs the sentence's value as $18$ numbers and the same value in the $\pm$ wording as $3.5$. The stated tolerance has little effect (Figure~\ref{fig:llm}). From $\pm 30$ to $\pm 10$ the Bayes-optimal strength rises ninefold (eightfold or more for the hedging predictor). The models' strength rises by a median factor of $1.3$, and by at most $1.6$ in $17$ of the $18$ pairs of model and wording; the rise is significant in $14$ of them (paired bootstrap). From $\pm 10$ to $\pm 3$, where the Bayes-optimal strength rises another elevenfold (at least threefold for the hedging predictor), the models' strength rises significantly in one pair and falls significantly in three. The models therefore respond to stated precision, but weakly and only at the loose end. They over-weigh a loose description in $14$ of the $18$ pairs, by up to $13$ numbers against $1$, and under-weigh a specific one, with at most $17$ numbers against $100$ ($16$ to $51$ for the hedging predictor). An exact line gets a similar weight ($0.7$ to $3.9$, against $1.0$ to $3.4$ at $\pm 3$). The strength also varies with $n$, where it is constant for the Bayes-optimal predictor that takes the line as valid: in $16$ of the $18$ pairs it is lower at $n = 41$ to $80$ than in our window, often after a rise at $n = 11$ to $20$; for the 32B instruction-tuned model's sentence it falls from $18$ to $6$. The response to the stated tolerance is as weak in later windows: at $n = 11$ to $20$ and at $n = 21$ to $40$ the median factor from $\pm 30$ to $\pm 10$ is $1.5$, and from $\pm 10$ to $\pm 3$ it is $1.1$ and $1.0$.

\paragraph{An invalid line is discounted, not discarded.} The models give the stated value of an invalid line $12\%$ to $71\%$ of the weight they give a valid one (Table~\ref{tab:llm}). The hedging Bayes-optimal predictor gives an invalid line a strength below $0.1$ in the same window, whether it starts from a prior probability of validity of $0.5$ or $0.9$; the 32B instruction-tuned model weighs an invalid sentence as $6$ numbers. Little of this weight is literal copying of the stated number: removing the probability that the predictive distribution puts on exactly that number lowers the strength of an invalid line by under $5\%$ in $17$ of the $18$ pairs and by $20\%$ in the other.

\paragraph{By regret, a description is worth little, and often less than nothing.} Under Definition~\ref{def:ess}, the ESS of a valid line at $\pm 10$ against the prompt with no line ranges from $-3.7$ to $4.2$ and is negative in $11$ of the $18$ pairs (Table~\ref{tab:llm}). A line changes what a model expects next, whatever it says. Directly after ``Mean: 508 $\pm$ 10'' the 7B base model puts $57\%$ of its probability on ``Median'': it reads the line as the first row of a table of summary statistics. With five numbers in hand and a valid line at $\pm 10$, the models put $9\%$ to $86\%$ of their probability on a number coming next, against $88\%$ to $92\%$ with no line. The content ESS, which holds the form of the prompt fixed, is $-0.1$ to $5.4$ numbers at $\pm 10$ ($-0.6$ to $3.6$ with the full next-token loss), where the hedging Bayes-optimal predictor's is $7$ to $11$. It also varies with $n$. In five of the six pairs of the 13B and 32B base models it grows with the numbers in hand, to $19$ for the 32B model's sentence at $n = 41$ to $80$, where the regret curves are nearly flat and a gap of a hundredth of a nat spans many numbers; in the other pairs it is negative, near zero, or undefined there.

\paragraph{The models learn slowly from numbers.} With no line, the regret after five numbers is $0.48$ to $1.30$ nats, against $0.086$ for the Bayes-optimal predictor, and after fifty it is $0.07$ to $0.47$ against $0.010$, where the instruction-tuned models are the slower ones. This agrees with \citet{liu2025}, who find that language models estimate a density from numbers in context like a kernel estimator of adaptive width. Slow learning from numbers raises the ESS of a description, which is an exchange rate, and the ESS is small even so.

\paragraph{Summary.} By the weight they give a stated value, the pretrained models treat a description as a number of demonstrations that is set by the model, by how the description is worded, and by how many demonstrations are in hand, and that barely tracks the precision the description states. The networks of \S\ref{sec:networks} were never told a description's precision; they learned it from the task distribution and came close to the Bayes-optimal ESS for reliable descriptions. The two experiments differ in more than training: the networks see a five-dimensional task and a description token, the language models a one-dimensional task in text, from one model family, with nothing that says how reliable a description is. With that caveat, the results are consistent with a description's worth to a learner being set by the learner's experience of how descriptions relate to tasks, more than by what a description says about its own precision.

\section{Conclusion}
We asked how many demonstrations a task description is worth, and answered for the Bayes-optimal predictor in in-context linear regression. A reliable description is worth a number of demonstrations that grows with the task dimension and with the description's specificity. A description that may be invalid is worth as much when it is valid and nothing when it is not, and not knowing which costs at most the entropy of its reliability; a predictor that does not hedge can lose more than the description gives. Networks trained on the task distribution come close to these values for reliable descriptions and take their trust in a description from that distribution. Pretrained language models that are told a description's precision in words respond to it only weakly: the weight they give a description is set mostly by the model and the wording.

\paragraph{Limitations.} The analysis assumes a linear-Gaussian task, a description that is equally specific in every direction, and a predictor that knows the description's specificity and reliability; a real description can be exact about some aspects of a task and silent about others. The Bayes-optimal numerical results are at one setting, $d = 16$ and $a_0 = 10$. The full-trust formula~\eqref{eq:trust} is for a prompt with no demonstrations, and Corollary~\ref{cor:ess}(ii) is a limit: with few demonstrations in hand, the ESS of a description that may be invalid can exceed $p$ times that of a reliable one. The network results rest on one training run per setting, two for one of them, and a network with a mixture output would not be bound by the single-Gaussian limit. The language-model test covers one model family and the simplest case of the setting; it scores plain-text continuation without a chat template, which is not how the instruction-tuned models are normally used; and nothing in its prompts says how reliable a description is. Both the implied strength and the regret-based ESS depend on the wording and on the number of demonstrations in hand; the implied strength is read from the mean of a predictive distribution that is wider than the Bayes-optimal one; and the regret-based ESS depends on the criterion: under squared error of the mean prediction, the content ESS at $\pm 10$ is $2.4$ to $23$ numbers, where under log loss it is $-0.1$ to $5.4$. The ESS is defined by the log-loss regret of the next prediction; another criterion would give a different number, and the bound on the price of not knowing whether a description is valid is specific to log loss.

\bibliographystyle{plainnat}
\bibliography{refs}

\clearpage
\appendix
\section{Proofs}

\subsection{Proof of Theorem \ref{prop:reliable} (Reliable Descriptions)}\label{app:specificity}
We state the theorem in full, with the regret bounds behind (i) and (ii), prove those bounds as two lemmas with their expansions, and locate the root that defines the ESS.

\medskip\noindent\textbf{Theorem~\ref{prop:reliable} (full statement).} {\itshape Let $p = 1$, let $\psi$ be the digamma function, and let $c = 1/(r a_0) = \sigma_y^2 / s_1^2$ with $s_1^2 = r s_0^2$: the ESS of the description's prior $\mathcal{N}(m, s_1^2 I_d)$ in the conjugate sense, noise variance over prior variance \citep{neuenschwander2020}, the conjugate-prior ESS.
\begin{enumerate}
\item[(i)] \emph{Loose descriptions.} Suppose $c \le 1$. Then for every $n < d$,
$R^{\mathrm{ex}}(n) \ge \tfrac12\big[\log(2a_0) + \psi\big(\tfrac{d-n}{2}\big)\big]$, with equality in the limit $a_0 \to \infty$, and as $rd \to \infty$, uniformly in $c \le 1$,
\[
\ess = (d-1)(1-r) + (1-r)\,c + O\!\left(\tfrac{1}{rd}\right).
\]
\item[(ii)] \emph{Specific descriptions.} Suppose $c \ge d$. Then for every $n \ge d$,
$R^{\mathrm{ex}}(n) \le \Phi(n) := \tfrac12\big[\psi\big(\tfrac{n+1}{2}\big) - \psi\big(\tfrac{n-d+1}{2}\big)\big]$, with equality in the limit $a_0 \to \infty$, and as $c/d \to \infty$, uniformly in $a_0 \ge 1$,
\[
\ess = c + d + 1 - \tfrac{1}{a_0} + O\!\left(\tfrac{d}{c}\right).
\]
\item[(iii)] For every $d$, $a_0$, and $r$, $\ess \le c + d + 2$.
\end{enumerate}}
\medskip

\paragraph{Standard facts.}
Throughout, $\sigma_y = 1$, $\varepsilon = 1/a_0$, $A = \|x\|^2 \sim \chi^2_d$ for the query $x$, $X_n \in \R^{n \times d}$ holds the demonstration inputs as rows, $\Sigma = (\varepsilon I + X_n^\top X_n)^{-1}$ is the posterior covariance given $n$ demonstrations and no description, and $g(t) = \E \log(1 + At)$, which is increasing and concave with
\begin{equation}\label{eq:g-derivs}
d\,(1 - (d+2)t) \le g'(t) = \E \frac{A}{1 + At} \le d,
\qquad
-(d^2 + 2d) \le g''(t) \le 0 .
\end{equation}
We use (a) $\E \log \chi^2_k = \psi(k/2) + \log 2$, $\E[\chi_k^{-2j}] = \prod_{i=1}^{j} (k - 2i)^{-1}$ for $k > 2j$, and $\psi(z) = \log z - 1/(2z) + O(z^{-2})$; (b) for $W \sim W_j(\nu, I)$ with $\nu \ge j$ and $z \sim \mathcal{N}(0, I_j)$ independent of $W$, $z^\top W^{-1} z$ is distributed as $\chi^2_j / \chi^2_{\nu-j+1}$ with independent numerator and denominator \citep[Theorem~3.2.12]{muirhead1982}, and $\E \operatorname{tr} W^{-2} = j(\nu-1)/((\nu-j)(\nu-j-1)(\nu-j-3))$ for $\nu - j \ge 4$ \citep{vonrosen1988}; (c) $y - y^2/2 \le \log(1 + y) \le y$ for $y \ge 0$.

\paragraph{Regret.}
The posterior predictive is the true conditional law of $y$, a Gaussian with variance $1 + x^\top \Sigma x$, so its expected log loss is its entropy and $R^{\mathrm{ex}}(n) = \tfrac12 \E \log(1 + x^\top \Sigma x)$, as in \S\ref{sec:specificity}. It is strictly decreasing in $n$, since by the Sherman--Morrison formula each new demonstration lowers $x^\top \Sigma x$ almost surely. For the description, $n = 0$ and the posterior covariance is $r a_0 I = I/c$, so by (a) and (c) with $y = c/A$,
\begin{equation}\label{eq:rdesc}
2R^{\mathrm{desc}} = g(1/c) = \log(2 r a_0) + \psi(d/2) + \delta^{\mathrm{desc}},
\qquad
\delta^{\mathrm{desc}} = \E \log(1 + c/A) = \frac{c}{d-2} + O\!\left(\frac{c^2}{d^2}\right) \quad (d \ge 5).
\end{equation}

\paragraph{Locating the root.}
Let $h(n) = 2R^{\mathrm{ex}}(n) - 2R^{\mathrm{desc}}$ on integers $n \ge 0$, with linear interpolant $\bar h$. Then $\ess$ is the unique root of $\bar h$, and $h(0) = g(a_0) - g(r a_0) > 0$. In (i) and (ii) we show, for some $n_0 \ge 0$, $\alpha > 0$, and $\theta \to 0$ along the stated limit,
\begin{equation}\label{eq:root}
h(n) = \alpha\,(n_0 - n) + O(\alpha\theta)
\qquad \text{uniformly over integers $n \ge 0$ with $|n - n_0| \le 3$.}
\end{equation}
Being a convex combination of two such values, $\bar h(t)$ obeys the same expansion at real $t \ge 0$ with $|t - n_0| \le 2$; so for small $\theta$ the root lies within $2$ of $n_0$, and $0 = \bar h(\ess) = \alpha(n_0 - \ess) + O(\alpha\theta)$ gives $\ess = n_0 + O(\theta)$.

\begin{lemma}[fewer than $d$ demonstrations]\label{lem:few}
Let $n < d$ and $k = d - n$. Then $2R^{\mathrm{ex}}(n) = \log(2a_0) + \psi(k/2) + \delta^{\mathrm{ex}}$ with $\delta^{\mathrm{ex}} \ge 0$, $\delta^{\mathrm{ex}} \to 0$ as $a_0 \to \infty$, and, for $k \ge 5$,
\[
\delta^{\mathrm{ex}} = \frac{\varepsilon (d-1)}{(k-1)(k-2)} + O\!\left(\frac{\varepsilon^2 d^2}{k^4}\right).
\]
\end{lemma}

\begin{proof}
Write $X_n^\top X_n = \sum_{i \le n} \lambda_i u_i u_i^\top$ with $\lambda_i > 0$ almost surely, and let $Q$ be the squared norm of the projection of $x$ onto the $k$ directions orthogonal to the $u_i$. Then $\Sigma$ has eigenvalue $a_0$ on those directions and $1/(\varepsilon + \lambda_i)$ on $u_i$, so
\[
1 + x^\top \Sigma x = a_0 Q + T, \qquad T = 1 + V, \quad V = \sum_{i \le n} \frac{z_i^2}{\varepsilon + \lambda_i}, \quad z_i = u_i^\top x,
\]
where, given $X_n$, the $z_i$ are i.i.d.\ standard normal and independent of $Q \sim \chi^2_k$. Taking logarithms and expectations, by (a),
\[
2R^{\mathrm{ex}}(n) = \log(2a_0) + \psi(k/2) + \delta^{\mathrm{ex}}, \qquad \delta^{\mathrm{ex}} = \E \log\!\left(1 + \frac{\varepsilon T}{Q}\right) \ge 0 .
\]
As $a_0 \to \infty$, $\varepsilon T$ decreases to $0$ pathwise, so $\delta^{\mathrm{ex}} \to 0$ by dominated convergence (the integrand at $a_0 = 1$ is integrable).

For the expansion, let $F = \sum_i z_i^2 / \lambda_i$. In the eigenbasis of $W = X_n X_n^\top \sim W_n(d, I)$, which has the same eigenvalues, $F = z^\top W^{-1} z$ and $\sum_i z_i^2/\lambda_i^2 = z^\top W^{-2} z$ with $z \sim \mathcal{N}(0, I_n)$ independent of $W$. By (b), $F \sim \chi^2_n / \chi^2_{k+1}$, so $\E F = n/(k-1)$ and $\E (1+F)^2 = O(d^2/k^2)$, and $\E \sum_i z_i^2/\lambda_i^2 = n(d-1)/(k(k-1)(k-3))$. Since $1/\lambda - \varepsilon/\lambda^2 \le 1/(\varepsilon + \lambda) \le 1/\lambda$, we have $F - \varepsilon \sum_i z_i^2/\lambda_i^2 \le V \le F$. Now (c) with $y = \varepsilon T / Q$, the independence of $Q$ from $T$, and (a) give
\[
\frac{\varepsilon\,(1 + \E V)}{k-2} - \frac{\varepsilon^2\, \E T^2}{2(k-2)(k-4)} \le \delta^{\mathrm{ex}} \le \frac{\varepsilon\,(1 + \E F)}{k-2} = \frac{\varepsilon (d-1)}{(k-1)(k-2)} ,
\]
and the two sides differ by $\varepsilon^2 \E \sum_i z_i^2/\lambda_i^2 /(k-2) + \varepsilon^2 \E T^2 / (2(k-2)(k-4)) = O(\varepsilon^2 d^2 / k^4)$, as $\E T^2 \le \E (1 + F)^2$.
\end{proof}

\begin{lemma}[at least $d$ demonstrations]\label{lem:many}
Let $n \ge d$ and $\ell = n - d + 1$. Then $2\Phi(n) = \E \log(1 + A/Q)$ with $Q \sim \chi^2_\ell$ independent of $A$, and $R^{\mathrm{ex}}(n) \le \Phi(n)$, with equality as $a_0 \to \infty$. If moreover $a_0 \ge 1$ and $\ell \ge \max(d, 7)$, then
\[
2R^{\mathrm{ex}}(n) = 2\Phi(n) - \frac{\varepsilon d}{\ell^2} + O\!\left(\frac{\varepsilon d^2}{\ell^3}\right).
\]
\end{lemma}

\begin{proof}
Rotate coordinates so that $x = \|x\| e_1$; the rows of $X_n$ remain i.i.d.\ $\mathcal{N}(0, I_d)$, independent of $A$. Then $x^\top \Sigma x = A\,\Sigma_{11} = A/S$ with $S$ the Schur complement of the first coordinate in $\Sigma^{-1}$, which the Woodbury identity writes, with $X_1$ the first column of $X_n$ and $X_{-1}$ the other $d-1$ columns, as
\[
S = \varepsilon + X_1^\top \big(I_n + a_0 X_{-1} X_{-1}^\top\big)^{-1} X_1 .
\]
Given $X_{-1}$, the middle matrix has eigenvalue $1$ on the $\ell$ directions orthogonal to the columns of $X_{-1}$ and $1/(1 + a_0 \mu_j)$ on the other $d - 1$, where $\mu_j$ are the eigenvalues of $M = X_{-1}^\top X_{-1} \sim W_{d-1}(n, I)$. Since $X_1 \sim \mathcal{N}(0, I_n)$ is independent of $X_{-1}$, its coordinates in that eigenbasis are i.i.d.\ standard normal, and
\[
S = Q + \Delta, \qquad Q \sim \chi^2_\ell, \qquad
\Delta = \varepsilon + \sum_{j < d} \frac{z_j^2}{1 + a_0 \mu_j} \in \big[\varepsilon,\, \varepsilon(1 + F)\big], \qquad F = \sum_{j < d} \frac{z_j^2}{\mu_j},
\]
with $A$, $Q$, and $(z, M)$ independent; by (b), $F \sim \chi^2_{d-1} / \chi^2_{\ell+1}$, so $\E F = (d-1)/(\ell-1)$ and, when $\ell \ge \max(d, 7)$, $\E (1+F)^2 \le 5$.

Since $A + Q \sim \chi^2_{n+1}$, (a) gives $\E \log(1 + A/Q) = \psi\big(\tfrac{n+1}{2}\big) - \psi\big(\tfrac{\ell}{2}\big) = 2\Phi(n)$. As $S \ge Q$, $R^{\mathrm{ex}}(n) \le \Phi(n)$, with equality as $a_0 \to \infty$ by monotone convergence, since $\Delta$ decreases to $0$.

For the expansion, $2R^{\mathrm{ex}}(n) = 2\Phi(n) - \E D$ with
\[
D = \log\Big(1 + \frac{A}{Q}\Big) - \log\Big(1 + \frac{A}{S}\Big) = \log(1 + y_1) - \log(1 + y_2), \quad y_1 = \frac{\Delta}{Q} \ge y_2 = \frac{\Delta}{Q + A}, \quad y_1 - y_2 = \frac{\Delta A}{Q(Q+A)} .
\]
By concavity of $\log(1+y)$ and $1/(1+y_1) \ge 1 - y_1$, $(y_1 - y_2)(1 - \Delta/Q) \le D \le y_1 - y_2$. Here $\Delta$ is independent of $(A, Q)$; $\E \frac{A}{Q(Q+A)} = \E\frac1Q - \E\frac{1}{Q+A} = \frac{d}{(\ell-2)(\ell+d-2)}$ by (a), as $Q + A \sim \chi^2_{\ell+d}$; $\frac{A}{Q^2(Q+A)} \le \frac{A}{Q^3}$; and $\varepsilon \le \E\Delta \le \varepsilon(1 + \E F) = \varepsilon\frac{\ell+d-2}{\ell-1}$, $\E \Delta^2 \le 5\varepsilon^2$. Hence
\[
\frac{\varepsilon d}{(\ell-2)(\ell+d-2)} - \frac{5\,\varepsilon^2 d}{(\ell-2)(\ell-4)(\ell-6)} \le \E D \le \frac{\varepsilon d}{(\ell-1)(\ell-2)} ,
\]
and both sides are $\varepsilon d / \ell^2 + O(\varepsilon d^2 / \ell^3)$, using $\varepsilon \le 1$.
\end{proof}

\paragraph{Any prior precision.}
The proof of Lemma~\ref{lem:many} uses $\varepsilon$ only as the precision of the prior. For $\tau > 0$ let $R^{\tau}(n) = \tfrac12 \E \log\big(1 + x^\top (\tau I + X_n^\top X_n)^{-1} x\big)$, the regret with $n$ demonstrations under a Gaussian prior of precision $\tau$ in every direction, so that $R^{\mathrm{ex}} = R^{\varepsilon}$. With $\tau$ in place of $\varepsilon$ and $1/\tau$ in place of $a_0$, the same proof gives $\Delta \in [\tau, \tau(1+F)]$ and, for $\ell \ge \max(d, 7)$,
\begin{equation}\label{eq:any-tau}
2R^{\tau}(n) = 2\Phi(n) - \frac{\tau d}{\ell^2} + O\!\left(\frac{(\tau + \tau^2)\, d^2}{\ell^3}\right).
\end{equation}

\begin{proof}[Proof of Theorem~\ref{prop:reliable}]
\emph{(iii).} Conditionally on $A$, $\log(1 + Av)$ is concave in $v$, so Jensen's inequality with $\E[1/Q] = 1/(\ell-2)$ gives $2\Phi(n) \le g(1/(\ell-2))$ for $\ell \ge 3$. If $\ell - 2 \ge c$, that is $n \ge c + d + 1$, then $g(1/(\ell-2)) \le g(1/c) = 2R^{\mathrm{desc}}$, so $R^{\mathrm{ex}}(n) \le \Phi(n) \le R^{\mathrm{desc}}$ by Lemma~\ref{lem:many}. The least such integer is $\lceil c + d + 1 \rceil < c + d + 2$, and $R^{\mathrm{ex}}$ is decreasing, so $\ess \le c + d + 2$.

\emph{(i).} The lower bound and the limit are Lemma~\ref{lem:few}. For the expansion let $c \le 1$, $n_0 = (d-1)(1-r) + (1-r)c$, and $n \ge 0$ an integer with $|n - n_0| \le 3$; write $k = d - n = rd + u$. Since $n_0 = d - rd - u_0$ with $u_0 = 1 - r + \varepsilon - c \in [-c, 1]$, we have $|u| \le 4$, so $k \ge 5$ once $rd \ge 9$. By Lemma~\ref{lem:few}, \eqref{eq:rdesc}, $\psi(z) = \log z - 1/(2z) + O(z^{-2})$, and $\varepsilon = rc$, with every $O$ uniform in $c \le 1$ and $|u| \le 4$,
\begin{align*}
h(n) &= \psi(k/2) - \psi(d/2) - \log r + \delta^{\mathrm{ex}} - \delta^{\mathrm{desc}} \\
&= \Big[\log\frac{k}{rd} - \frac1k + \frac1d\Big] + \frac{\varepsilon d}{r^2 d^2}\Big(1 + O\Big(\frac{1}{rd}\Big)\Big) - \frac{c}{d} + O\Big(\frac{1}{(rd)^2}\Big)
 = \frac{u - 1 + r + c - \varepsilon}{rd} + O\Big(\frac{1}{(rd)^2}\Big),
\end{align*}
since $\log(k/(rd)) = u/(rd) + O((rd)^{-2})$, $1/k = 1/(rd) + O((rd)^{-2})$, $\varepsilon d / (r^2 d^2) = c/(rd)$, and $c/d = \varepsilon/(rd)$. So $h(n) = \frac{1}{rd}\big[(n_0 - n) + O(1/(rd))\big]$, which is~\eqref{eq:root} with $\alpha = \theta = 1/(rd)$, and $\ess = (d-1)(1-r) + (1-r)c + O(1/(rd))$.

\emph{(ii).} The upper bound and the limit are Lemma~\ref{lem:many}. For the expansion let $n_0 = c + d + 1 - \varepsilon$ and $n$ an integer with $|n - n_0| \le 3$, so that $\ell - 2 = c + s$ with $s = n - n_0 - \varepsilon$, $|s| \le 4$; for $c \ge 9d$ every such $n$ has $\ell \ge \max(d, 7)$. By~\eqref{eq:rdesc} and Lemma~\ref{lem:many}, $2\Phi(n) - 2R^{\mathrm{desc}} = \E\big[g(1/Q) - g(1/c)\big]$, and a second-order Taylor expansion of $g$ about $1/c$ with~\eqref{eq:g-derivs}, $\E[1/Q] - 1/c = -s/(c(c+s))$, and $\E(1/Q - 1/c)^2 = \operatorname{Var}(1/Q) + s^2/(c(c+s))^2 = O(c^{-3})$ gives
\[
2\Phi(n) - 2R^{\mathrm{desc}} = -\frac{s}{c(c+s)}\,d\,\Big(1 + O\Big(\frac dc\Big)\Big) + O\Big(\frac{d^2}{c^3}\Big) = -\frac{ds}{c^2} + O\Big(\frac{d^2}{c^3}\Big).
\]
With Lemma~\ref{lem:many}, $\varepsilon \le 1$, and $\varepsilon d/\ell^2 = \varepsilon d/c^2 + O(d/c^3)$,
\[
h(n) = -\frac{d}{c^2}\Big[s + \varepsilon + O\Big(\frac{d}{c}\Big)\Big] = \frac{d}{c^2}\Big[(n_0 - n) + O\Big(\frac{d}{c}\Big)\Big],
\]
which is~\eqref{eq:root} with $\alpha = d/c^2$ and $\theta = d/c$, uniformly in $a_0 \ge 1$, so $\ess = c + d + 1 - 1/a_0 + O(d/c)$.
\end{proof}

\subsection{Proof of Proposition \ref{prop:unreliable} (Descriptions That May Be Invalid)}\label{app:reliability}
Let $D_n = (x_{1:n}, y_{1:n}, m, x)$ be the prompt with the query input, and let $\pi_n = P(Z = 1 \mid D_n)$ be the posterior probability that the description is valid. Write $\pi_n(Z)$ for the posterior weight on the true case: $\pi_n$ if $Z = 1$ and $1 - \pi_n$ otherwise.

\paragraph{Step 1: the mixture predictive.}
The Bayes-optimal predictive is $f(y) = \pi_n f_1(y) + (1 - \pi_n) f_0(y)$, where $f_1$ and $f_0$ are the predictives of the two cases given $D_n$: $f_1$ is the predictive of the Bayes-optimal predictor for a reliable description, and $f_0$ is that of the Bayes-optimal predictor with the demonstrations alone, because under $Z = 0$ the description is independent of $w$ and carries no information about $y$. At $n = 0$, $\pi_0 = p$: seeing $m$ says nothing about $Z$, because $m$ has the same law $\mathcal{N}(0, (1-r) a_0 I)$ in both cases, and $x$ is independent of $Z$. The two predictives are then $f_1 = \mathcal{N}(m^\top x,\, 1 + r a_0\|x\|^2)$ and $f_0 = \mathcal{N}(0,\, 1 + a_0\|x\|^2)$.

\paragraph{Step 2: the decomposition.}
Consider a predictor that is told $Z$. It predicts with $f_1$ when $Z=1$, incurring regret $R^{\mathrm{desc}}_{p=1}(n)$, and with $f_0$ when $Z=0$, incurring regret $R^{\mathrm{ex}}(n)$, so its expected regret is $p R^{\mathrm{desc}}_{p=1}(n) + (1-p) R^{\mathrm{ex}}(n)$. Given $(D_n, Z)$ the true conditional density of $y$ is $f_Z$, so the Bayes-optimal predictor's expected log loss exceeds the told predictor's by
\[
\E\big[\log f_Z(y) - \log f(y)\big] = \E\big[\mathrm{KL}(f_Z \,\|\, f)\big] = I(Z;\, y \mid D_n) = I_n \ge 0,
\]
the conditional mutual information, because $f$ is the law of $y$ given $D_n$ and $f_Z$ its law given $D_n$ and $Z$. This is~\eqref{eq:decomp}.

\paragraph{Step 3: the bounds.}
By definition $I_n = H(Z \mid D_n) - H(Z \mid D_n, y)$. A query input is independent of $Z$ and of everything else in the prompt, so $H(Z \mid D_n) = H(Z \mid x_{1:n}, y_{1:n}, m)$, and the query of $D_n$ with its label is a further demonstration, so $H(Z \mid D_n, y) = H(Z \mid D_{n+1})$. Hence $I_n = H(Z \mid D_n) - H(Z \mid D_{n+1}) \le H(Z \mid D_n)$, the entropies $H(Z \mid D_n)$ are non-increasing in $n$, and the sum telescopes: $\sum_{n < N} I_n = H(Z \mid D_0) - H(Z \mid D_N) \le H(Z \mid D_0) = H(p)$, as $\pi_0 = p$ by Step 1.

\paragraph{Step 4: case by case.}
Let the label of each query join the demonstrations, so that $D_{n+1}$ extends $D_n$. By Bayes' rule, $f_Z(y)/f(y) = \pi_{n+1}(Z)/\pi_n(Z)$: the excess log loss on a label is the log of the factor by which that label raises the posterior weight on the true case. Summed over the first $N$ labels the excess telescopes to $\log \pi_N(Z) - \log \pi_0(Z) \le -\log \pi_0(Z)$, and $\pi_0 = p$ by Step 1. So over the whole prompt the Bayes-optimal predictor's log loss exceeds the told predictor's by at most $\log(1/p)$ when the description is valid and by at most $\log(1/(1-p))$ when it is invalid.

This completes the proof of Proposition~\ref{prop:unreliable}.

\begin{proof}[Proof of Corollary~\ref{cor:ess}]
\emph{(i).} At $n = 0$, \eqref{eq:decomp} with $I_0 \ge 0$ gives $R^{\mathrm{desc}}(0) \ge (1-p)R^{\mathrm{ex}}(0)$ for every $r$, because $R^{\mathrm{desc}}_{p=1}(0) = \tfrac12 \E \log(1 + r a_0 \|x\|^2)$ is nonnegative. Since $R^{\mathrm{ex}}(n)$ is strictly decreasing in $n$ and tends to $0$ (Lemma~\ref{lem:many}, as $\Phi(n) \to 0$), there is a finite $n^*_p$ with $R^{\mathrm{ex}}(n^*_p) = (1-p)R^{\mathrm{ex}}(0)$, and the root that defines $\ess_0$ cannot exceed it. Moreover $R^{\mathrm{desc}}_{p=1}(0)$ decreases to $0$ as $r \to 0$ by monotone convergence, so in that limit the standalone regret lies between $(1-p)R^{\mathrm{ex}}(0)$ and $(1-p)R^{\mathrm{ex}}(0) + H(p)$.

\emph{(ii).} Fix $d$, $a_0$, and $r$; every $O$ below is as $n \to \infty$ with these fixed. A reliable description leaves the prior $\mathcal{N}(m, r a_0 I)$, of precision $c$ in every direction, and the regret does not depend on $m$, so $R^{\mathrm{desc}}_{p=1} = R^{c}$ in the notation of~\eqref{eq:any-tau}, while $R^{\mathrm{ex}} = R^{\varepsilon}$. Since $1/\ell^2 = 1/n^2 + O(n^{-3})$, \eqref{eq:any-tau} gives
\[
R^{\mathrm{ex}}(n) - R^{\mathrm{desc}}_{p=1}(n) = \frac{(c - \varepsilon)\,d}{2n^2} + O(n^{-3}).
\]
By the standard expansion $\psi(z) = \log z - 1/(2z) - 1/(12z^2) + O(z^{-4})$, $\psi(z + \tfrac12) - \psi(z) = 1/(2z) + 1/(8z^2) + O(z^{-3})$, so $2\Phi(n) - 2\Phi(n+1) = d/(\ell(n+1)) + O(n^{-3})$, and with~\eqref{eq:any-tau} again,
\[
R^{\mathrm{ex}}(n) - R^{\mathrm{ex}}(n+1) = \frac{d}{2n^2} + O(n^{-3}).
\]
Let $L(n) = p\,R^{\mathrm{desc}}_{p=1}(n) + (1-p)\,R^{\mathrm{ex}}(n) = R^{\mathrm{ex}}(n) - p\,\big(R^{\mathrm{ex}}(n) - R^{\mathrm{desc}}_{p=1}(n)\big)$ be the first two terms of~\eqref{eq:decomp}, and let $n + k_n$ be the root at which the interpolated $R^{\mathrm{ex}}$ equals $L(n)$. On any bounded number of steps past $n$ the interpolated $R^{\mathrm{ex}}$ falls by $\tfrac{d}{2n^2}(1 + O(1/n))$ per step, and it must fall by $p\,(c - \varepsilon)\,\tfrac{d}{2n^2}(1 + O(1/n))$ in total, so $k_n = p\,(c - \varepsilon) + O(1/n)$. If $p = 1$, then $R^{\mathrm{desc}} = L$ and $\ess_n = k_n \to c - \varepsilon$. If $p < 1$, then $R^{\mathrm{desc}}(n) = L(n) + I_n \ge L(n)$ and $R^{\mathrm{ex}}$ is decreasing, so $\ess_n \le k_n$ and $\limsup_n \ess_n \le p\,(c - \varepsilon)$.
\end{proof}

\subsection{The Fully Trusting Predictor}\label{app:trust}
We derive~\eqref{eq:trust}. With no demonstrations, the fully trusting predictor uses $f_1 = \mathcal{N}(m^\top x, v_1)$ with $v_1 = 1 + r a_0 \|x\|^2$. If the description is valid, $f_1$ is the true conditional law of $y$ and the regret is $\tfrac12 \E \log v_1 = R^{\mathrm{desc}}_{p=1}(0)$. If it is invalid, $w$ is independent of $m$, so given $x$ the error $y - m^\top x$ has mean zero and variance $1 + a_0\|x\|^2 + (1-r)a_0\|x\|^2 = v_1 + 2(1-r)a_0\|x\|^2$. The expected log loss of $f_1$ is then $\tfrac12 \log(2\pi v_1) + \tfrac12 + (1-r)a_0\|x\|^2/v_1$, against $\tfrac12 \log(2\pi) + \tfrac12$ for the oracle, so the regret given $x$ is $\tfrac12 \log v_1 + (1-r)a_0\|x\|^2/v_1$. Averaging over the two cases gives~\eqref{eq:trust}. The integrand of the penalty increases to $a_0\|x\|^2$ as $r \to 0$, so the penalty increases to $(1-p)\,a_0 d$ by monotone convergence.

\section{The Language-Model Test at Every Tolerance}\label{app:llm}
Table~\ref{tab:llm-full} gives the implied strength and the content ESS of \S\ref{sec:llm} for every model, wording, and stated tolerance.

\begin{table}[t]
\centering\small
\setlength{\tabcolsep}{4pt}
\caption{Implied strength of a valid and an invalid line, and content ESS of a valid line over the number tokens and with the full next-token loss, for every model, wording, and stated tolerance ($2{,}000$ tasks, $n = 5, \dots, 10$). The Bayes-optimal rows are the predictor that takes the line as valid and the predictor that hedges with a prior probability of validity $p$.}
\label{tab:llm-full}
\vspace{3pt}
\input{../../results/llm-test/llm_table_full.tex}
\end{table}

\end{document}
